% ==========================================
%  style/chapters.tex  —  章节标题格式
%  基于 ctex 宏集与 titlesec 宏包
% ==========================================

% --- 确保章节从奇数页开始 ---
\newcommand{\ensureoddstart}{%
	\clearpage
	\ifodd\value{page}
	% 奇数页，什么也不做
	\else
		\hbox{}
		\newpage
	\fi
}

% --- 章节 + 摘要排版 ---
\newcommand{\customchapter}[2]{%
	\ensureoddstart
	\chapter{#1}%
	\vspace*{-\baselineskip}%
	\vspace*{4\baselineskip}%
	\indent #2%
	\vspace*{2\baselineskip}%
	\newpage
}

% --- 无序号节标题 ---
\newcommand{\subsec}[1]{%
	\refstepcounter{subsection}%
	\setcounter{subsubsection}{0}%
	\subsection*{\thesubsection~#1}%
}

\newcommand{\ssubsec}[1]{%
	\refstepcounter{subsubsection}%
	\subsubsection*{\thesubsubsection~#1}%
}

\newcommand{\tsubsec}[1]{\subsubsection*{\textreferencemark~#1}}

% --- 浮动体与章节标题的英文名称 ---
\addto\captionsenglish{
	\renewcommand{\contentsname}{\heiti{\zihao{-2} 目\hspace*{2em}录}}
	\renewcommand{\listfigurename}{插图目录}
	\renewcommand{\listtablename}{表格目录}
	\renewcommand{\refname}{参考文献}
	\renewcommand{\abstractname}{摘要}
	\renewcommand{\indexname}{索引}
	\renewcommand{\tablename}{表}
	\renewcommand{\figurename}{图}
}

% --- 正文标题层级格式 ---
\ctexset{
	chapter = {
		name = {,},
		number = {第\,\chinese{chapter}\,章},
		format = \kaishu\centering\zihao{-1},
		aftername = \par,
		beforeskip = 18pt,
		afterskip = 18pt
	},
	section = {
		name = {,},
		number = {\arabic{chapter}.\arabic{section}},
		format = \songti\raggedright\zihao{4},
		aftername = \hspace{1em},
		beforeskip = 12pt,
		afterskip = 6pt
	},
	subsection = {
		name = {,},
		number = {\arabic{chapter}.\arabic{section}.\arabic{subsection}},
		format = \heiti\raggedright\zihao{-4},
		aftername = \hspace{1em},
		beforeskip = 12pt,
		afterskip = 6pt
	},
}
